%%%%%%%%%%%%%%%%%%%%%%%%%%%%%%%%%%%%%%%%%%%%%%%%%%%%%%%%%%%%%%
%% Rozdział 2: Przegląd literatury
%%%%%%%%%%%%%%%%%%%%%%%%%%%%%%%%%%%%%%%%%%%%%%%%%%%%%%%%%%%%%%

\section{Przegląd literatury}
\label{chap:related-work}

\subsection{Rozpoznawanie i tłumaczenie języka migowego (SLR/SLT)}
\label{sec:slr-slt-background}

Automatyczne tłumaczenie języka migowego dzieli się na etapy z wykorzystaniem
pośredniej reprezentacji glosowej, wprowadzonej w rozdziale~\ref{chap:introduction}
(podrozdział~\ref{sec:motivation-context}). Podział ten prowadzi do czterech zadań
kierunkowych: \emph{wideo--glosy} i \emph{glosy--tekst} po stronie rozpoznawania
oraz \emph{tekst--glosy} i \emph{glosy--wideo} z wykorzystaniem awatara po stronie
generowania. Niniejszy projekt koncentruje się najpierw na etapie wideo--glosy
(rozdział~\ref{chap:isolated-gloss-recognition}), a następnie obejmuje tłumaczenie
między glosami a tekstem (rozdział~\ref{chap:text-translation}) oraz integrację
obu komponentów w potok przetwarzania na poziomie zdań
(rozdział~\ref{chap:sentence-level-translation}).

Rozpoznawanie glos na podstawie wideo jest nietrywialne nawet w wariancie
izolowanym, w którym każde nagranie wejściowe zawiera pojedynczy, uprzednio
wyodrębniony znak. Znak odpowiadający tej samej glosie jest wykonywany odmiennie
przez różne osoby migające: różnice w układzie dłoni, trajektorii ruchu, tempie
migania i proporcjach ciała powodują znaczną zmienność wewnątrzklasową.
Konwencjonalne klasyfikatory z funkcją softmax, które dzielą przestrzeń wejściową
za pomocą ustalonych granic decyzyjnych, słabo radzą sobie z uwzględnianiem tej
zmienności, szczególnie przy ograniczonej ilości danych, typowej dla korpusów
języka migowego. Obserwacja ta uzasadnia przyjęte w niniejszej pracy podejście do
rozpoznawania izolowanych glos oparte na wyszukiwaniu i uczeniu metryki
(zob. rozdział~\ref{chap:isolated-gloss-recognition}), w odróżnieniu od dominujących
podejść klasyfikacyjnych omówionych poniżej.

\subsection{Metody przekształcania wideo w glosy}
\label{sec:video-to-gloss-methods}

W tym podrozdziale przedstawiono przegląd proponowanych w literaturze architektur
rozpoznawania izolowanych znaków migowych oraz ciągłych wypowiedzi migowych,
obejmujący podejścia oparte na szkielecie, wyglądzie oraz rozwiązania hybrydowe.

W pracy \emph{Sign Pose-Based Transformer for Word-Level Sign Language Recognition}
\cite{bohacek2022spoter} przedstawiono SPOTER -- system, który całkowicie pomija
surowe klatki RGB i operuje wyłącznie na dwuwymiarowych punktach kluczowych
szkieletu, uzyskanych za pomocą estymatora pozy (stawy ciała i dłoni). Każda klatka
jest reprezentowana jako wektor współrzędnych stawów; sekwencja tych wektorów
z całego nagrania stanowi wejście standardowego transformera o architekturze
enkoder--dekoder. Pojedynczy uczony token zapytania o klasę wykorzystuje mechanizm
uwagi nad zakodowaną sekwencją klatek (podobnie jak w podejściu DETR do detekcji
obiektów), aby uzyskać końcową klasyfikację. Do kluczowych elementów rozwiązania
należą etap normalizacji, który przeskalowuje współrzędne stawów względem stałego
układu odniesienia ciała w celu wyeliminowania zmienności wynikającej z różnic
między osobami i ich odległości od kamery, oraz schemat augmentacji dostosowany
do danych szkieletowych (losowy obrót stawów, skalowanie i zniekształcenia zbliżone
do perspektywicznych, stosowane bezpośrednio do punktów kluczowych zamiast do
pikseli). Raportowane wyniki wskazują na dokładność top-1 wynoszącą 63,18\%
na większym zbiorze WLASL, obejmującym wiele osób migających, wobec 100\%
na mniejszym zbiorze LSA64, zarejestrowanym w bardziej kontrolowanych warunkach.
Pokazuje to, że lekka architektura pozwala znacząco ograniczyć rozmiar modelu
oraz koszt wnioskowania kosztem pewnego spadku dokładności na zróżnicowanych danych
o dużym zasobie słownictwa.

Praca \emph{Power Data Classification: A Hybrid of a Novel Local Time Warping and LSTM}
\cite{li2016ltwlstm} wykorzystuje standardowe podejście LSTM do klasyfikacji
sekwencji, z funkcją softmax stosowaną do sumy stanów ukrytych. Jej wyróżniającym
wkładem jest jednak Local Time Warping (LTW), autorska miara odległości, która
zastępuje dopasowanie DTW oparte na programowaniu dynamicznym nieprzemiennym
porównaniem lokalnym w oknie o stałej szerokości -- zmniejszając złożoność z $O(nm)$
do czasu liniowego przy zachowaniu porównywalnej zdolności rozróżniania klas.
Końcowy klasyfikator łączy na poziomie prawdopodobieństw wyniki klasyfikatora
jednego najbliższego sąsiada wykorzystującego LTW z wyjściem softmax sieci LSTM.
W pięciokrotnej walidacji krzyżowej samodzielny model LSTM osiąga dokładność około
85--88\%, porównywalną z samym klasyfikatorem najbliższego sąsiada opartym na LTW,
natomiast połączenie obu gałęzi zwiększa dokładność do 89--93\%. Autorzy przypisują
ten wynik strukturalnie odmiennym błędom popełnianym przez gałąź parametryczną
i nieparametryczną. Taki schemat łączenia dwóch gałęzi jest strukturalnie
analogiczny do połączenia wyuczonego enkodera czasowego z wyszukiwaniem opartym
na FAISS, stanowiącego jedno z rozszerzeń omówionych
w podrozdziale~\ref{sec:gloss-recognition-discussion}.

W pracy \emph{A Comprehensive Study on Deep Learning-Based Methods for Sign Language
Recognition} \cite{adaloglou2021comprehensive} ponownie zaimplementowano
 i wytrenowano kilka rodzin modeli w identycznych warunkach eksperymentalnych
(te same zbiory danych, przetwarzanie wstępne i podziały danych), aby oddzielić
wpływ architektury od czynników zakłócających, które zwykle różnią się między
publikacjami. Oceniane rodziny obejmują klasyfikatory klatek oparte na 2D-CNN,
sieci 3D-CNN (np.\ sploty czasoprzestrzenne w stylu I3D), hybrydy CNN+RNN oraz
modele grafowe operujące na punktach kluczowych pozy. Oceniono je osobno dla
rozpoznawania izolowanych znaków oraz rozpoznawania ciągłych wypowiedzi w języku
migowym (CSLR, zwykle trenowanego z funkcją straty CTC). Główny wniosek jest taki,
że żadna architektura nie dominuje we wszystkich zbiorach danych i zadaniach:
modele oparte na pozie i grafach zwykle lepiej uogólniają na różne osoby migające,
natomiast modele CNN/3D-CNN oparte na wyglądzie zazwyczaj osiągają wyższe wyniki
na zbiorach o ograniczonym zróżnicowaniu osób migających. Podkreśla to silną
zależność raportowanej dokładności od zbioru danych. Ten wniosek metodologiczny
bezpośrednio uzasadnia przeprowadzone w rozdziale~\ref{chap:isolated-gloss-recognition}
kontrolowane porównanie pięciu enkoderów czasowych w jednakowych warunkach.

Szerszy przegląd architektur obejmujący zarówno rozpoznawanie, jak i tłumaczenie
\cite{review2024dlprocessing} omawia podejścia 3D-CNN, w których czasoprzestrzenne
jądra splotowe stosuje się bezpośrednio do stosów klatek wideo, oraz systemy
tłumaczenia języka migowego oparte na transformerach (np.\ Camgoza i in.). Te
ostatnie zastępują potok CNN+RNN mechanizmem samouwagi nad cechami wideo lub
wyodrębnionymi sekwencjami glos, eliminując ograniczenie wynikające z sekwencyjnego
przetwarzania w sieciach RNN. Przegląd obejmuje również multimodalne modele
bazowe wykorzystujące uczenie transferowe, w których enkodery wizualne są
wstępnie trenowane na dużych korpusach rozpoznawania czynności (np.\ Kinetics),
a następnie dostrajane na mniejszych zbiorach danych języka migowego. Podejście
to konsekwentnie poprawia dokładność rozpoznawania w zadaniu docelowym względem
trenowania od podstaw.

Kilka kolejnych badań charakteryzuje praktyczne możliwości podejść opartych
na wyglądzie i rozwiązań hybrydowych. Potok GoogLeNet--BiLSTM trenowany na
nagraniach terenowych, przy niekontrolowanym oświetleniu i tle
\cite{agri2021bilstm}, osiąga średnią dokładność 76,21\% na autorskim zbiorze danych
z dziedziny rolnictwa. Jest to wynik wyraźnie niższy od raportowanych dla
porównywalnych modeli hybrydowych trenowanych na nagraniach studyjnych, co autorzy
przypisują niekontrolowanym warunkom pozyskiwania danych, a nie słabości
architektury. Porównanie klasyfikatorów CNN, MLP, RNN, SVM i KNN na statycznych
obrazach układów dłoni \cite{ann2022recognition} wykazuje skuteczność 99\%
na zbiorze treningowym liczącym 50\,000 obrazów, jednak odzwierciedla to dopasowanie
do danych treningowych, a nie zdolność uogólniania; ponadto zastosowany model
nie obsługuje ciągłych wypowiedzi ani znaków dynamicznych. Przegląd uwzględniający
modalności danych \cite{modalities2022review} wyróżnia systemy wizyjne (RGB/RGB-D
jako wejście potoków CNN lub CNN-RNN), systemy oparte na szkielecie lub pozie
(sekwencje punktów kluczowych jako wejście sieci GCN lub transformerów) oraz
systemy wykorzystujące czujniki noszone na ciele (IMU/EMG jako wejście sieci RNN).
Autorzy zauważają, że w benchmarkach multimodalnych, takich jak AUTSL, zazwyczaj
raportuje się wyższą dokładność niż na jednomodalnych zbiorach RGB dzięki
dodatkowej informacji rozróżniającej klasy, pochodzącej ze strumieni głębi
 i danych szkieletowych.

Dwie kolejne architektury bezpośrednio łączą etapy splotowe i rekurencyjne.
Moduł wejściowy 3D-CNN połączony z warstwami LSTM \cite{cnn3dlstm2023realtime}
wykonuje splot jednocześnie w wymiarach wysokości, szerokości i czasu, aby
wyodrębnić cechy przestrzenne uwzględniające ruch bez osobnego etapu obliczania
przepływu optycznego. Następnie przekazuje uzyskaną sekwencję cech dla danego
nagrania do warstw LSTM w celu modelowania zależności czasowych o większym zasięgu.
\emph{Deepsign} \cite{deepsign2022} stosuje pojedynczą warstwę LSTM, po której
następuje warstwa GRU, do przetwarzania cech wyznaczanych na poziomie klatek.
W badaniu empirycznym wykazano, że spośród czterech możliwych dwuwarstwowych
konfiguracji LSTM/GRU najlepsze wyniki daje układ LSTM, a następnie GRU, osiągający
około 97\% dokładności dla jedenastu klas gestów na autorskim zbiorze danych.

Wreszcie w pracy \emph{Spatial-Temporal Graph Convolutional Networks for Sign Language
Recognition} \cite{correia2019stgcn} pozycja szkieletu w każdej klatce jest
reprezentowana jako graf, którego węzłami są stawy, a krawędziami połączenia
anatomiczne (np.\ nadgarstek--łokieć, łokieć--bark). Przestrzenny splot grafowy
agreguje informacje między połączonymi stawami w obrębie klatki, natomiast splot
czasowy łączy ten sam staw w kolejnych klatkach, zgodnie z podejściem ST-GCN,
pierwotnie opracowanym do ogólnego rozpoznawania czynności na podstawie szkieletu.
W ten sposób struktura kinematyczna jest jawnie zakodowana jako założenie
indukcyjne, w przeciwieństwie do podejść, które otrzymują jedynie spłaszczony
wektor współrzędnych dla każdej klatki. Brak takiego jawnego kodowania stanowi
ograniczenie architektury modelu przyjętej
w rozdziale~\ref{chap:isolated-gloss-recognition}; możliwość jego wprowadzenia
ponownie omówiono w podrozdziale~\ref{sec:gloss-recognition-discussion} jako
kierunek dalszych prac.

\subsection{Modele NLP dla języka migowego}
\label{sec:nlp-models-sign-language}

Poza rozpoznawaniem wizualnym tłumaczenie glos na tekst i tekstu na glosy pozwala
ująć przetwarzanie języka migowego jako problem NLP typu sekwencja--sekwencja.
Zazwyczaj rozwiązuje się go za pomocą architektur enkoder--dekoder opartych
na transformerach, operujących na sekwencjach tokenów glos, a nie wyłącznie
na surowym tekście (zob. systemy tłumaczenia oparte na transformerach wspomniane
w podrozdziale~\ref{sec:video-to-gloss-methods}
\cite{review2024dlprocessing}).

\todo{DO UZUPEŁNIENIA: Należy przedstawić przegląd modeli NLP stosowanych
bezpośrednio do tłumaczenia glos na tekst i tekstu na glosy (gloss2text /
text2gloss) -- np. modeli tłumaczenia języka migowego typu sekwencja--sekwencja
oraz opartych na transformerach (Camgoz i in., Sign Language Transformers),
dostrajania wstępnie wytrenowanych modeli językowych do tłumaczenia glos,
tłumaczenia wstecznego oraz strategii augmentacji danych dla niewielkich korpusów
glosowych. Ten podrozdział powinien stanowić uzasadnienie wyborów dotyczących
modelowania dokonanych w rozdziale~\ref{chap:text-translation}.}

\subsection{Istniejące systemy i benchmarki}
\label{sec:existing-systems-benchmarks}

Omówione publikacje wykorzystują różnorodne zbiory danych, które stanowią również
punkt odniesienia przy projektowaniu porównania modeli w niniejszej pracy.
\textbf{WLASL} i \textbf{LSA64} wykorzystano w pracy \cite{bohacek2022spoter}
do porównania wyników uzyskiwanych dla dużego zasobu słownictwa prezentowanego
przez wiele osób migających z wynikami dla mniejszego zasobu zarejestrowanego
w kontrolowanych warunkach. \textbf{AUTSL} to duży multimodalny zbiór danych
tureckiego języka migowego, wykorzystywany do oceny modeli bazowych CNN-LSTM
 i 3D-CNN \cite{modalities2022review}. \textbf{ASLLVD} stanowi podstawę zbioru danych
szkieletowych udostępnionego wraz z modelem ST-GCN \cite{correia2019stgcn}.
\textbf{IISL2020} to autorski zbiór danych obejmujący jedenaście klas, używany
do oceny architektury z ułożonymi kolejno warstwami LSTM i GRU
w pracy \cite{deepsign2022}. W niniejszej pracy porównanie modeli rozpoznawania
izolowanych glos przeprowadzono na podzbiorze zbioru \textbf{ASL Citizen}
(Microsoft, 2023), szczegółowo opisanym
w podrozdziale~\ref{sec:gloss-recognition-dataset}.

\todo{DO UZUPEŁNIENIA: Należy opisać istniejące całościowe systemy SLR/SLT
oraz publiczne benchmarki i rankingi wyników (np. How2Sign,
RWTH-PHOENIX-Weather 2014T, istniejące systemy produkcyjne lub demonstracyjne
systemy badawcze), wykraczając poza omówione już zbiory danych, oraz odnieść je
do zakresu niniejszej pracy.}